\newpage
\section{Gliederung}

Die nachfolgende Grobgliederung orientiert sich am sechsphasigen Prozessmodell der Design Science Research Methodology \autocite[Vgl.][S.~45]{Peffers2007} und folgt dem in der gestaltungsorientierten Wirtschaftsinformatik etablierten Aufbau empirisch-konstruktiver Arbeiten \autocite[Vgl.][S.~75]{Hevner2004}.

{\renewcommand{\labelenumii}{\theenumi.\arabic{enumii}}
\begin{enumerate}
  \item Einleitung
    \begin{enumerate}
      \item Problemstellung und Motivation
      \item Zielsetzung und Forschungsfrage
      \item Aufbau der Arbeit
    \end{enumerate}

  \item Theoretische Grundlagen
    \begin{enumerate}
      \item Robotic Process Automation: Begriff, Konzept und Abgrenzung
      \item UI-basierte RPA: Technologie und Funktionsweise
      \item Eignung von Geschäftsprozessen für RPA
      \item API-lose Legacy-Systeme: Charakteristika und Integrationsherausforderungen
    \end{enumerate}

  \item Forschungsstand und verwandte Arbeiten
    \begin{enumerate}
      \item Bestehende Vorgehensmodelle zur RPA-Implementierung
      \item RPA in Legacy-Umgebungen: bisherige Ansätze und Lücken
      \item Identifikation des Forschungsdesiderats
    \end{enumerate}

  \item Forschungsmethodik
    \begin{enumerate}
      \item Design Science Research als methodischer Rahmen
      \item Artefakttyp und Evaluierungsansatz
      \item Vorgehen bei der Modellentwicklung und -validierung
    \end{enumerate}

  \item Entwicklung des Vorgehensmodells
    \begin{enumerate}
      \item Anforderungsanalyse an ein kontextspezifisches Vorgehensmodell
      \item Phasenstruktur des Modells
      \item Tätigkeiten, Rollen und Artefakte je Phase
      \item Besonderheiten des UI-RPA-Einsatzes in API-losen Umgebungen
    \end{enumerate}

  \item Evaluation des Vorgehensmodells
    \begin{enumerate}
      \item Evaluierungsstrategie
      \item Durchführung der Expertenvalidierung oder Fallstudie
      \item Ergebnisse und Implikationen
    \end{enumerate}

  \item Diskussion
    \begin{enumerate}
      \item Beantwortung der Forschungsfrage
      \item Grenzen der Arbeit und Übertragbarkeit
      \item Implikationen für Forschung und Praxis
    \end{enumerate}

  \item Fazit und Ausblick
\end{enumerate}
}

% Leitfragen (FOM-Gliederungshilfe):
% - Hier nur eine GROBGLIEDERUNG, die Ihren momentanen Wissensstand darstellt.
% - Noch kein zu differenziertes Inhaltsverzeichnis. Die Gliederung verfeinert
%   sich im Abstimmungsprozess und entwickelt sich waehrend der Bearbeitung weiter.

% TODO: Grobgliederung der geplanten Thesis (z.B. als Aufzaehlung).
% \begin{enumerate}
%   \item Einleitung
%   \item ...
%   \item Fazit
% \end{enumerate}
